\section{Reduced Hilbert-scheme DT/PT/OP comparison}
\label{sec:reduced-dt-pt-op-comparison}

\subsection{Claim attacked}

The claim under attack is the comparison
\begin{equation}
\label{eq:rdtpt-claim}
Z^X_{\mathrm{HilbDT}}=Z^X_{\mathrm{OP}}
\end{equation}
in the quotient-by-\(E\), Behrend-weighted normalisation.  The attack
has a positive outcome: \eqref{eq:rdtpt-claim} follows from
Oberdieck's reduced stable-pair theorem for \(S\times E\), together
with the primitive Gromov--Witten/pairs comparison used by
Oberdieck--Pixton.

\subsection{Variables and objects}

Let
\begin{equation}
\label{eq:rdtpt-x-class}
X=S\times E,\qquad
(\beta_h,d)=\iota_{S*}\beta_h+\iota_{E*}(d[E]),
\qquad
\langle\beta_h,\beta_h\rangle=2h-2,
\end{equation}
where \(\beta_h\neq0\) is primitive.  Use the OP/Siegel variables
\begin{equation}
\label{eq:rdtpt-op-vars}
P=e^{iu}=e^{2\pi iz},\qquad
Q=e^{2\pi i\tau},\qquad
T=e^{2\pi i\widetilde\tau}.
\end{equation}
The Borcherds variables of this manuscript are
\begin{equation}
\label{eq:rdtpt-borch-vars}
q_{\mathrm Bch}=Q,\qquad r_{\mathrm Bch}=P,\qquad s_{\mathrm Bch}=T.
\end{equation}
The manuscript's displayed DT variables
\[
Z^X_{\mathrm{DT}}(q,t,p)
=
\sum_{h,d,n}\mathbf{DT}^X_{n,(\beta_h,d)}
q^{d-1}t^{h-1}(-p)^n
\]
therefore satisfy
\begin{equation}
\label{eq:rdtpt-manuscript-map}
q=T,\qquad t=Q,\qquad p=P.
\end{equation}

\subsection{Reduced Hilbert-scheme DT}

Bryan defines the ordinary ideal-sheaf invariant by
\begin{equation}
\label{eq:rdtpt-bryan-ordinary}
DT_{\beta,n}(X)
=
\int_{\operatorname{Hilb}^{\beta,n}(X)}\nu\,de,
\end{equation}
and observes that for \(X=S\times E\) the ordinary invariant vanishes
because the \(E\)-translation action on
\(\operatorname{Hilb}^{\beta+dE,n}(X)\) is fixed-point free.  The reduced
Hilbert-scheme invariant is therefore
\begin{equation}
\label{eq:rdtpt-bryan-reduced}
\mathsf{DT}^{\mathrm{Hilb}}_{\beta+dE,n}(X)
=
\int_{\operatorname{Hilb}^{\beta+dE,n}(X)/E}\nu\,de.
\end{equation}
Bryan's partition function is
\begin{equation}
\label{eq:rdtpt-bryan-series}
\mathsf{DT}^{\mathrm{Bryan}}(p_{\mathrm{Br}},q_{\mathrm{Br}},
\widetilde q_{\mathrm{Br}})
=
\sum_{h,d,n}
\mathsf{DT}^{\mathrm{Hilb}}_{h,d,n}(X)\,
\widetilde q_{\mathrm{Br}}^{\,h-1}
q_{\mathrm{Br}}^{d-1}(-p_{\mathrm{Br}})^n.
\end{equation}
Bryan states that his \(q_{\mathrm{Br}},\widetilde q_{\mathrm{Br}}\)
are opposite to Oberdieck--Pandharipande's variables.  Thus
\begin{equation}
\label{eq:rdtpt-bryan-map}
p_{\mathrm{Br}}=P,\qquad
q_{\mathrm{Br}}=T,\qquad
\widetilde q_{\mathrm{Br}}=Q.
\end{equation}
Equivalently,
\begin{equation}
\label{eq:rdtpt-hilbdt-series}
Z^X_{\mathrm{HilbDT}}(T,Q,P)
:=
\sum_{h,d,n}
\mathsf{DT}^{\mathrm{Hilb}}_{h,d,n}(X)\,
T^{d-1}Q^{h-1}(-P)^n
=
\mathsf{DT}^{\mathrm{Bryan}}(P,T,Q).
\end{equation}

\subsection{Reduced stable pairs}

Let \(P_n(X,(\beta,d))\) be the moduli space of stable pairs
\((F,s)\) on \(X\) with
\[
\chi(F)=n,\qquad [\operatorname{Supp}(F)]=(\beta,d).
\]
Oberdieck defines the quotient reduced stable-pairs invariant by
\begin{equation}
\label{eq:rdtpt-pt-quotient}
\mathsf P^X_{n,(\beta,d)}
=
\int_{P_n(X,(\beta,d))/E}\nu\,de,
\end{equation}
where \(de\) is orbifold topological Euler characteristic and \(\nu\)
is the Behrend function on the quotient stack.  Oberdieck proves that
this quotient invariant agrees with the incidence invariant obtained
from the reduced virtual class:
\begin{equation}
\label{eq:rdtpt-pt-incidence}
\mathsf P^X_{n,(\beta,d)}
=
\int_{[P_n(X,(\beta,d))]^{\mathrm{red}}}
\tau_0\!\left(\pi_1^*(\beta^\vee)\cup\pi_2^*[\mathrm{pt}]\right),
\qquad
\langle\beta,\beta^\vee\rangle=1.
\end{equation}
This is Oberdieck, \emph{On reduced stable pair invariants},
arXiv:1605.04631, Theorem~1.  The proof constructs a symmetric perfect
obstruction theory on an isotrivial fibre \(K=p^{-1}(0_E)\), pushes the
resulting virtual class to \(P_n(X,(\beta,d))/E\), and compares it with
the reduced virtual class by the Fulton--Siebert formula.

\subsection{Reduced DT/PT theorem}

For every non-zero class \(\beta\in H_2(S,\mathbb Z)\) and every
\(d\geq0\), Oberdieck defines
\begin{equation}
\label{eq:rdtpt-fixed-series}
\mathsf{DT}_{(\beta,d)}(x)
=
\sum_n\mathsf{DT}^{\mathrm{Hilb}}_{\beta+dE,n}(X)x^n,
\qquad
\mathsf{PT}_{(\beta,d)}(x)
=
\sum_n\mathsf P^X_{n,(\beta,d)}x^n.
\end{equation}
Oberdieck, arXiv:1605.04631, Theorem~3(i), proves
\begin{equation}
\label{eq:rdtpt-oberdieck-dtpt}
\mathsf{DT}_{(\beta,d)}(x)=\mathsf{PT}_{(\beta,d)}(x).
\end{equation}
The same theorem records the degree-zero issue: in the local wall
crossing one first has
\[
\mathsf{DT}_C(x)=\mathsf{PT}_C(x)\,\mathsf{DT}_0(x),
\qquad
\mathsf{DT}_0(x)=M(-x)^{\chi(X)}.
\]
For \(X=S\times E\),
\[
\chi(X)=\chi(S)\chi(E)=24\cdot0=0,
\qquad
\mathsf{DT}_0(x)=1.
\]
Hence \eqref{eq:rdtpt-oberdieck-dtpt} has no residual MacMahon factor.

Substituting \(x=-P\) in \eqref{eq:rdtpt-oberdieck-dtpt} gives, for
each primitive \(\beta_h\) and \(d\geq0\),
\begin{equation}
\label{eq:rdtpt-coefficient-comparison}
\sum_n
\mathsf{DT}^{\mathrm{Hilb}}_{h,d,n}(X)(-P)^n
=
\sum_n
\mathsf P^X_{n,h,d}(-P)^n.
\end{equation}
Multiplying by \(Q^{h-1}T^{d-1}\) and summing over \(h,d\) gives
\begin{equation}
\label{eq:rdtpt-hilbdt-pt-series}
Z^X_{\mathrm{HilbDT}}(T,Q,P)
=
Z^X_{\mathrm{PT,red}}(P,Q,T),
\end{equation}
where
\begin{equation}
\label{eq:rdtpt-pt-series}
Z^X_{\mathrm{PT,red}}(P,Q,T)
:=
\sum_{h,d,n}
\mathsf P^X_{n,h,d}\,Q^{h-1}T^{d-1}(-P)^n.
\end{equation}

\subsection{PT/OP bridge}

Oberdieck--Pandharipande define the reduced stable-pairs series
\[
\mathsf P^X_{\beta,d}(y)=\sum_n\mathsf P^X_{n,\beta,d}y^n
\]
and prove, for primitive \(\beta_h\) and all \(d\), the reduced
Gromov--Witten/pairs correspondence
\begin{equation}
\label{eq:rdtpt-gw-pt}
\sum_n\mathsf P^X_{n,\beta_h,d}y^n
=
\sum_g\mathsf N^{X\bullet}_{g,h,d}u^{2g-2},
\qquad
y=-e^{iu}=-P.
\end{equation}
This is the proposition immediately following Conjecture~D in
Oberdieck--Pandharipande, arXiv:1411.1514.  Oberdieck--Pixton use the
same quotient stable-pairs convention:
\begin{equation}
\label{eq:rdtpt-op-pt}
\mathsf P_{n,(\beta,d)}
=
\int_{P_n(X,(\beta,d))/E}\nu\,de,
\qquad
\sum_n\mathsf P_{n,h,d}y^n
=
\sum_g\mathsf N_{g,h,d}u^{2g-2}.
\end{equation}
Therefore
\begin{equation}
\label{eq:rdtpt-pt-op}
Z^X_{\mathrm{PT,red}}(P,Q,T)
=
\mathcal Z_{\mathrm{OP}}(u,Q,T).
\end{equation}
Oberdieck--Pixton Theorem~1 gives
\begin{equation}
\label{eq:rdtpt-op-theorem}
\mathcal Z_{\mathrm{OP}}(u,Q,T)
=
-\frac{1}{\chi_{10}^{\mathrm{OP}}(P,Q,T)}.
\end{equation}
With the manuscript normalisation
\begin{equation}
\label{eq:rdtpt-chi-normalisation}
\chi_{10}^{\mathrm{OP}}(P,Q,T)
=D_5(Z)^2
=4096^{-1}\Delta_5(Z)^2,
\qquad
(q_{\mathrm Bch},r_{\mathrm Bch},s_{\mathrm Bch})=(Q,P,T),
\end{equation}
one obtains
\begin{equation}
\label{eq:rdtpt-op-constant}
\mathcal Z_{\mathrm{OP}}(u,Q,T)
=-4096\,\Delta_5(Z)^{-2}.
\end{equation}

\subsection{Comparison theorem}

\begin{theorem}[Reduced Hilbert-scheme DT/PT/OP comparison]
\label{thm:rdtpt-comparison}
For \(X=S\times E\), primitive non-zero \(\beta_h\), \(d\geq0\), and
the variables \eqref{eq:rdtpt-op-vars}--\eqref{eq:rdtpt-manuscript-map},
the quotient-by-\(E\), Behrend-weighted Hilbert-scheme DT partition
function equals the Oberdieck--Pixton primitive reduced partition
function:
\begin{equation}
\label{eq:rdtpt-main}
Z^X_{\mathrm{HilbDT}}(T,Q,P)
=
Z^X_{\mathrm{PT,red}}(P,Q,T)
=
\mathcal Z_{\mathrm{OP}}(u,Q,T).
\end{equation}
Consequently
\begin{equation}
\label{eq:rdtpt-main-constant}
Z^X_{\mathrm{HilbDT}}(T,Q,P)
=
-4096\,\Delta_5(Z)^{-2},
\qquad
C_{\mathrm{DT}}=C_{\mathrm{OP}}=-4096.
\end{equation}
\end{theorem}

\begin{proof}
Equation \eqref{eq:rdtpt-hilbdt-pt-series} is Oberdieck's reduced
DT/PT theorem after the substitution \(x=-P\).  Equation
\eqref{eq:rdtpt-pt-op} is the primitive reduced
Gromov--Witten/pairs correspondence used in Oberdieck--Pixton.  Equation
\eqref{eq:rdtpt-op-constant} is Oberdieck--Pixton Theorem~1 after
\(\chi_{10}^{\mathrm{OP}}=4096^{-1}\Delta_5^2\).  Combining these three
identities proves \eqref{eq:rdtpt-main} and
\eqref{eq:rdtpt-main-constant}.  The sign \((-P)^n\) is not an
additional convention: it is the common substitution \(y=-e^{iu}=-P\)
in \eqref{eq:rdtpt-gw-pt} and the variable used in Bryan's reduced DT
series \eqref{eq:rdtpt-bryan-series}.  The degree-zero factor is
\(M(-P)^{\chi(S\times E)}=1\).
\end{proof}

\subsection{Replacement text for the manuscript}

Replace the Hilbert-scheme DT scalar-normalisation assumption by the
following theorem.

\begin{theorem}[Reduced Hilbert-scheme DT scalar square]
\label{thm:hilb-dt-op-square-replacement}
Let \(X=S\times E\), and let
\(\beta_h\in H_2(S,\mathbb Z)\) be primitive and non-zero with
\(\langle\beta_h,\beta_h\rangle=2h-2\).  Define
\[
\mathbf{DT}^X_{n,(\beta_h,d)}
=
\int_{\operatorname{Hilb}^n(X,(\beta_h,d))/E}\nu\,de
\]
and
\[
Z^X_{\mathrm{DT}}(q,t,p)
=
\sum_{h,d,n}\mathbf{DT}^X_{n,(\beta_h,d)}
q^{d-1}t^{h-1}(-p)^n.
\]
Under the variable map
\[
p=P=e^{iu},\qquad t=Q,\qquad q=T,
\qquad
(q_{\mathrm Bch},r_{\mathrm Bch},s_{\mathrm Bch})=(Q,P,T),
\]
one has
\[
Z^X_{\mathrm{DT}}(T,Q,P)
=
\mathcal Z_{\mathrm{OP}}(u,Q,T)
=
-4096\,\Delta_5(Z)^{-2}.
\]
Equivalently,
\[
C_{\mathrm{DT}}=-4096.
\]
\end{theorem}

\begin{proof}
Oberdieck's reduced DT/PT theorem for \(S\times E\)
(arXiv:1605.04631, Theorem~3(i)) gives
\[
\sum_n\mathbf{DT}^X_{n,(\beta_h,d)}(-P)^n
=
\sum_n\mathsf P^X_{n,(\beta_h,d)}(-P)^n.
\]
The usual degree-zero DT factor is
\[
M(-P)^{\chi(S\times E)}=M(-P)^0=1.
\]
Oberdieck--Pandharipande's primitive reduced
Gromov--Witten/pairs comparison gives
\[
\sum_n\mathsf P^X_{n,(\beta_h,d)}(-P)^n
=
\sum_g\mathsf N^X_{g,h,d}u^{2g-2}.
\]
After multiplying by \(Q^{h-1}T^{d-1}\) and summing in \(h,d\), this is
the OP partition function.  Oberdieck--Pixton Theorem~1 gives
\[
\mathcal Z_{\mathrm{OP}}(u,Q,T)
=
-\left(\chi_{10}^{\mathrm{OP}}(P,Q,T)\right)^{-1}.
\]
Finally,
\[
\chi_{10}^{\mathrm{OP}}(P,Q,T)=D_5(Z)^2=4096^{-1}\Delta_5(Z)^2.
\]
\end{proof}

Replace the status table entries by:
\[
\begin{array}{c|c|c}
\hbox{Object} & \hbox{Equation} & \hbox{Status} \\
\hline
\hbox{Primitive reduced curve counting}
& Z^X_{\mathrm{OP}}=-4096\Delta_5^{-2}
& \hbox{Oberdieck--Pixton theorem} \\
\hbox{Hilbert-scheme DT series}
& Z^X_{\mathrm{DT}}(T,Q,P)=Z^X_{\mathrm{OP}}(P,Q,T)
& \hbox{Oberdieck reduced DT/PT plus OP primitive GW/PT} \\
\hbox{Scalar-square constant}
& C_{\mathrm{DT}}=C_{\mathrm{OP}}=-4096
& \hbox{theorem; no MacMahon factor since }\chi(S\times E)=0 \\
\hline
\end{array}
\]

Replace every occurrence of
\[
Z^X_{\mathrm{DT}}=\frac{C_{\mathrm{DT}}}{\Delta_5^2},
\qquad C_{\mathrm{DT}}\in\mathbb C^\times
\]
by
\[
Z^X_{\mathrm{DT}}(T,Q,P)
=
-4096\,\Delta_5(Z)^{-2}.
\]
Replace every occurrence of ``the Hilbert-scheme DT branch is
conditional'' by ``the Hilbert-scheme DT branch is theorem-level after
Oberdieck's reduced DT/PT theorem for \(S\times E\).''

\subsection{Source anchors}

\[
\begin{array}{c|c|c}
\hbox{source} & \hbox{anchor} & \hbox{use} \\
\hline
\hbox{Bryan, arXiv:1504.02920}
& \hbox{Definition 2.1}
& \eqref{eq:rdtpt-bryan-reduced},\ \eqref{eq:rdtpt-bryan-series} \\
\hbox{Oberdieck, arXiv:1605.04631}
& \hbox{Theorem 1}
& \eqref{eq:rdtpt-pt-incidence} \\
\hbox{Oberdieck, arXiv:1605.04631}
& \hbox{Theorem 3(i)}
& \eqref{eq:rdtpt-oberdieck-dtpt},\ \mathsf{DT}_0=1 \\
\hbox{Oberdieck--Pandharipande, arXiv:1411.1514}
& \hbox{Proposition after Conjecture D}
& \eqref{eq:rdtpt-gw-pt} \\
\hbox{Oberdieck--Pixton, arXiv:1706.10100}
& \hbox{Theorem 1 and equation (3)}
& \eqref{eq:rdtpt-op-theorem},\ \eqref{eq:rdtpt-chi-normalisation}
\end{array}
\]

\subsection{Remaining obligations}

The mathematical comparison gap is closed subject to the hypotheses
explicit in the cited theorems: \(\beta_h\) is non-zero and primitive,
\(d\geq0\), and the reduced quotient invariants use the Behrend
function of the quotient stack.  The integration obligations are:
\begin{enumerate}
\item add a bibliography entry for Oberdieck, \emph{On reduced stable
pair invariants}, arXiv:1605.04631;
\item replace the scalar-square assumption in the manuscript by
Theorem~\ref{thm:hilb-dt-op-square-replacement};
\item keep the protected-index dictionary separate: the enumerative
identity \(Z^X_{\mathrm{DT}}=Z^X_{\mathrm{OP}}\) is proved, while the
identification with a microscopic D-brane Hilbert space remains
physical.
\end{enumerate}
